\documentclass[convert]{standalone}
\usepackage[RPvoltages]{circuitikz}
\begin{document}
\begin{circuitikz}
\ctikzset{resistors/scale=0.8}

% Nombre neutro: el diagrama no identifica el modo de operacion.
\draw (0,0) node[op amp, noinv input up](OA){};

% Primera fuente y divisor conectado a la entrada no inversora.
\draw (-7,0) node[ground]{}
to[american voltage source, l_=$V_a$] (-7,2)
to[R=$R_1$, -*] (-4,2) coordinate(P)
to[R, l_=$R_2$] (-4,0.5) node[ground]{};
\draw (P) -- (-1.8,2) |- (OA.+);

% Segunda fuente y divisor conectado a la entrada inversora.
\draw (-7,-4) node[ground]{}
to[american voltage source, l_=$V_b$] (-7,-2)
to[R=$R_3$, -*] (-4,-2) coordinate(M)
to[R, l_=$R_4$] (-4,-3.5) node[ground]{};
\draw (M) -- (-1.8,-2) |- (OA.-);

% Alimentacion simetrica y salida sin conexion a las entradas.
\draw (OA.up) -- ++(0,0.5) node[vcc]{$+5\ V$};
\draw (OA.down) -- ++(0,-0.5) node[vee]{$-5\ V$};
\draw (OA.out) to[short, -o] ++(1.5,0) node[above]{$v_o$};

\end{circuitikz}
\end{document}
